\chapter{Preface}

This book has a simple ambition: a single, coherent book of Earth and space
sciences covering everything from the first year of school to the end of
the bachelor's degree, written in English for readers anywhere in the
world.

The style is concise and rigorous: courses built from definitions, examples,
propositions, methods and diagrams, with careful reasoning whenever it is
accessible at the given level, followed by carefully chosen exercises.
Results stated without full derivation are explicitly marked as admitted.
Full solutions to all exercises are collected at the end of the book.
Every number quoted about the real Earth, its oceans, its atmosphere or
the sky comes from a reference source, and every curve drawn from real
measurements is computed from the published data.

The course comes in five books. The first covers the primary and middle
school years, the second the three years of high school, and the last
three the three years of university. The two school books never teach the
same notion twice: each idea is taught once, in the year where it first
fits, and a later chapter that needs it opens with a short reminder,
\emph{You already know}, before going further. Where a notion belongs to
physics, biology or chemistry, the companion books of the same collection
teach it, and this book recalls it before taking the view of the Earth
scientist. In the school books each year is a part, split into chapters;
each university book is one year, split into chapters. The earlier the
year, the younger the reader it is written for: more illustrations, more
detailed worked steps, and gentler exercises.

\section*{How to read this book}

Statements are numbered within each chapter and color-coded: definitions in
blue, theorems in red, propositions and their relatives in orange, and
methods --- short recipes for standard tasks --- in green. Exercises are
graded from ($\star$) for direct applications to ($\star\star\star$) for
harder problems. Framed boxes hold what is not part of the argument: a
reminder of what an earlier chapter taught, how geologists and other Earth
scientists work in the field or in the laboratory, a page of history,
figures that change from year to year (with the year they are given for),
and safety advice.

\section*{Safety}

The Earth is learnt by looking at it: the sky, the weather, a landscape, the
stones of a building, the soil of a garden. That is all this book ever asks
of its reader. Never look at the Sun, with or without an instrument: it
blinds. Quarries, cliffs, caves, rivers in flood and volcanoes are places
for professionals with training and equipment, and this book describes their
work without sending anyone there.

\section*{Contributing}

The source of this book lives in a public git repository:
\begin{center}
\url{https://github.com/bvirrion/one-earth-and-space-sciences-book}
\end{center}
Contributions --- new chapters, corrections, better explanations,
additional exercises --- are welcome; see \texttt{CONTRIBUTING.md} at the
root of the repository.
